%% ============================================================
%% 分片 16：原文第 6 节「Loop and Line Operators」（原文 glshort.tex 起于第 4483 行）
%% 由 tools/harvmac2tex.py 从 arXiv 源码（hep-th/0604151v3）自动生成。
%% 只译英文散文；公式、编号、\label、\cite 一律不改。
%% ============================================================

\subsection{Topological 't Hooft Operators}

\label{sec:tophooft}

We have found Wilson operators in the topological field theory at
$t=\pm i$, that is at $\Psi=\infty$.  By  the $S$-duality
$S:\Psi\to -1/n_{\frak g}\Psi$, there must be dual operators at
$\Psi=0$. The topological Wilson operators were classified by the
choice of a representation $R$ of $G$.  So, as $S$ exchanges $G$
and $^L\negthinspace G$, the dual operators that appear in the
TQFT at $\Psi=0$ must be classified by a representation
$^L\negthinspace R$ of $^L\negthinspace G$.

The Wilson operator is an example of what is often called an
``order'' operator in statistical mechanics.  An order operator is
built from a classical expression (such as the holonomy of the
connection around a loop) that is then interpreted as a quantum
operator and included as a factor in a path integral.

The dual of an order operator is frequently a ``disorder''
operator.  A disorder operator is defined by specifying a
singularity that the classical fields are supposed to have, and
performing a path integral (or quantizing) in the presence of the
singularity.  Classically, the order and disorder operators seem
like completely different kinds of things.  But quantum
mechanically, they turn out to have, in suitable cases, the same
formal properties, and are frequently exchanged by duality. In our
problem, the dual of a Wilson operator is an 't Hooft  operator,
which is an example of a disorder operator.\footnote{For an explicit
demonstration that the dual of a Wilson operator is an 't Hooft
operator in the case of $G=U(1)$, see \cite{ref89}.}

To define a Wilson loop operator associated with a loop ${\cal
S}\subset M$, we required an orientation of ${\cal S}$.  The 't
Hooft loop operator instead requires an orientation of the normal
bundle to ${\cal S}$.  In fact, it is convenient to identify a
small neighborhood of ${\cal S}$ with ${\cal S}\times \Bbb{R}^3$
and to write $\epsilon_3$ for a volume-form on $\Bbb{R}^3$.  The
definition of the 't Hooft operator will depend only on the
orientation of the normal bundle, not the details of these
choices. If $r$ denotes the distance from ${\cal S}$ (in some
metric), then we can also write % 式 (6.6)
\setcounter{equation}{5}
\begin{equation}\label{eq:congo}
\epsilon_3=dr\,d{\rm
Vol},
\end{equation}
 where $d{\rm Vol}$ is a volume-form on the two-sphere at
fixed $r$.

Throughout this paper, when we write explicit formulas for 't Hooft operators in $U(N)$
or $SU(N)$ gauge theory, we consider the gauge fields as connections on a rank $N$ vector
bundle.  This means that $A$ and $\phi$ are are one-forms valued in
$N\times N$ skew-hermitian matrices; for $N=1$, they are simply imaginary one-forms.

Let us first describe the most basic 't Hooft operator in $U(1)$
gauge theory. On $\Bbb{R}^4=\Bbb{R}\times \Bbb{R}^3$, we pick
coordinates $(s,\vec x)$.  The singular Dirac monopole is a
classical solution of Maxwell's equations that is invariant under
translations of $s$ and rotations of $\vec x$, and has a
singularity at the line $L$ defined by $\vec x=0$. If $F=dA$ is
the curvature of a $U(1)$ connection $A$, then Maxwell's equations
read $dF=d\star F=0$, and can be solved on $\Bbb{R}^3\backslash
\{0\}$ (that is, on the complement of the point $\vec x=0$ in
$\Bbb{R}^3$) by % 式 (6.7)
\setcounter{equation}{6}
\begin{equation}\label{eq:izo}
F={i\over 2}\star d\left({1\over |\vec
x|}\right).
\end{equation}
  The use of the Hodge $\star$ operator shows
explicitly the dependence on the orientation of the normal bundle
to $L$.  To get a solution on $\Bbb{R}^4\backslash L$, we simply
use the projection $\Bbb{R}^4\backslash L\to \Bbb{R}^3\backslash
\{0\}$, and pull $F$ back to a solution of Maxwell's equations on
$\Bbb{R}^4\backslash L$.
 $F$ has been normalized so that
if $Y$ is a two-sphere enclosing the singularity at $\vec x=0$,
then $\int_Y \,iF/2\pi =1$.  Hence there exists a unitary line
bundle ${\cal L}\to \Bbb{R}^4\backslash L$  that admits a
connection of curvature $F$. This is the singular Dirac monopole.
It is characterized by $\int_Yc_1({\cal L})=1$.

The basic 't Hooft operator of abelian gauge theory is defined by
saying that the path integral is to be performed over connections
on a line bundle over $\Bbb{R}^4\backslash L$ that agree with the
singular Dirac monopole near $L$. Here there is no need to
consider only straight lines. $L$ can be replaced by any closed
curve  ${{\cal S}}$ in a general four-manifold $M$, or (if $M$ is
geodesically complete) any curve that goes off to infinity. The
path integral with an insertion of the 't Hooft loop operator
$T_{(0)}$ on the curve ${{\cal S}}$ is obtained by performing a
path integral for abelian gauge fields\footnote{If the homology class
of ${\cal S}$ in $H_1(M,\Bbb{Z})$ is nonzero, there are no such
abelian gauge fields and we declare that this path integral
vanishes. The dual of this is that if the homology class of ${\cal
S}$ is nonzero, the Wilson operator $W(R,{\cal S})$ has a
vanishing expectation value after integrating over all
connections; this vanishing follows from the behavior in twisting
by a flat line bundle.} on $M\backslash {{\cal S}}$ such that the
curvature near ${{\cal S}}$ diverges as $F\sim {i\over 2}\star_3d
(1/r)$, where $r=|\vec x|$ is now the distance from ${{\cal S}}$
and $\star_3$ is an operation on forms that near ${{\cal S}}$
looks like the $\star$ operator on planes normal to ${{\cal S}}$.
More generally, picking an integer $m$ and requiring that the
singular part of $F$ is $F\sim{(im/ 2)}\star_3 d (1/r)$, we get
the 't Hooft loop operator $T_{(0)}(m)$ of charge $m$.

Now let us consider ${\EUN}=4$ super Yang-Mills theory with gauge
group $U(1)$.  We can define an 't Hooft operator by asking that
the gauge field have the Dirac monopole singularity while other
fields are non-singular.  But is such an operator compatible with
the topological symmetry at $\Psi=0$?

We can decide this question as follows. To get to $\Psi=0$, we
take $t=1$ and $\tau$ imaginary. In order for the 't Hooft
operator to preserve the topological symmetry, the singular
behavior of the fields near ${{\cal S}}$ must be compatible with
that symmetry.  Since the question concerns the local behavior
near ${{\cal S}}$, it is convenient to temporarily revert to the
case that ${{\cal S}}$ is a line $L\subset \Bbb{R}^4$.

 The conditions for a set of fields
$(A,\phi)$ to preserve the topological symmetry were found in
eqns. \eqref{eq:loopme}.  If those equations are compatible with the type of
singularity required by the 't Hooft operator, then the 't Hooft
operator preserves supersymmetry. Otherwise, it is impossible for
any field configuration in the presence of the 't Hooft operator
to be
 supersymmetric, and we say that the 't Hooft operator violates
supersymmetry.

At $t=1$ and for abelian gauge theory, the requisite conditions
read % 式 (6.8)
\setcounter{equation}{7}
\begin{equation}\label{eq:nurot}
(F+d\phi)^+=(F-d\phi)^-=0.
\end{equation}
 These equations are
not satisfied if $F$ has the Dirac monopole singularity and $\phi$
is non-singular.  To get an 't Hooft operator that preserves the
topological symmetry, just as in the Wilson case, we have to
include $\phi$ in a suitable fashion.

What sort of singularity of $\phi$ must accompany the Dirac
monopole singularity of $F$ in order for the equations \eqref{eq:nurot}\ to
be satisfied? This question has a simple answer: we should take
$\phi=(i/2|\vec x|)ds$, where $s$ is the ``time'' coordinate along
$L$.  Thus, the equations \eqref{eq:nurot}\ are obeyed if we take
% 式 (6.9)
\setcounter{equation}{8}
\begin{equation}\label{eq:hildo}
\begin{aligned}F& ={i\over 2}\star_3 d{1\over |\vec x|}\\
                    \phi& = {i\over 2|\vec x|} ds,\\\end{aligned}
\end{equation}

where $\star_3$ is the three-dimensional Hodge $\star$ operator in
the directions normal to $L$.

This tells us how to define, for a general closed one-manifold
${{\cal S}}\subset M$, the basic 't Hooft operator $T$ that
preserves the topological symmetry at $\Psi=0$: to evaluate a
matrix element of this operator, we perform a path integral over
fields that possess near ${{\cal S}}$ the singularity described in
\eqref{eq:hildo}. Similarly, for any integer $m$, to get a charge $m$
topological 't Hooft operator $T(m)$, we simply require that the
fields have a singularity along ${{\cal S}}$ that is $m$ times
what is described in \eqref{eq:hildo}.

\par\medskip{\it Generalization To Any $G$ And Langlands
Duality}

The next step is to generalize this construction to arbitrary $G$.
This is done in a simple fashion \cite{ref90}.  We
simply pick an arbitrary homomorphism $\rho:U(1)\to G$.   Then we
define a topological 't Hooft operator by asking that the fields
should have a singularity along ${{\cal S}}$ that looks like the
image under $\rho$ of the singularity in \eqref{eq:hildo}.

The operator defined this way depends on $\rho$ only up to
conjugacy.  But conjugacy classes of homomorphisms $\rho:U(1)\to
G$ are classified by highest weights of the dual group
$^L\negthinspace G$.  In other words, they are classified by
irreducible representations of $^L\negthinspace G$.

So for every irreducible representation $^L\negthinspace R$ of
$^L\negthinspace G$, and every one-manifold ${{\cal S}}$, we have
obtained a topological 't Hooft operator $T({}^L\negthinspace
R,{{\cal S}})$. As in the Wilson case, we refer to this operator
as just $T({}^L\negthinspace R)$ or $T$ if the context is clear.
The fact that 't Hooft operators of $G$ are classified by
representations of $^L\negthinspace G$ is the classic GNO duality
\cite{ref10}, as recently reinterpreted in terms of operators rather than
states \cite{ref90}.

The classification of 't Hooft operators is also intimately
related to Grothendieck's classification of $G$-bundles on ${\Bbb
{CP}}^1$.  According to this classification, a $G$-bundle on
${\Bbb {CP}}^1$ is associated to the fundamental line bundle
${\cal O}(1)\to\Bbb{CP}^1$ via a homomorphism $\rho:U(1)\to G$ (or
its complexification $\rho:\Bbb{C}^*\to G_{\Bbb{C}}$).  Let us
identify with $\Bbb{CP}^1$ a small two-sphere surrounding the
singularity that defines an 't Hooft operator. Every $G$-bundle on
a Riemann surface, even if unstable, admits a connection obeying
the Yang-Mills equations \cite{ref91}, and on $\Bbb{CP}^1$ every
solution of the Yang-Mills equations is equivalent to one of the
abelian solutions used to construct 't Hooft operators.  So the
classification of 't Hooft operators is the same as the
classification of $G$-bundles on $\Bbb{CP}^1$.

\par\medskip{\it Topologically Non-Trivial 't Hooft
Operators}

A holomorphic $G$-bundle over $S^2=\Bbb{CP}^1$ may be
topologically non-trivial, though the topological classification
of $G$-bundles is much less fine than the holomorphic
classification.

A $G$-bundle over $S^2$ is classified by a characteristic class
$\xi\in H^2(S^2,\pi_1(G))\cong \pi_1(G)$.  But $\pi_1(G)$ is
naturally isomorphic to the center of $^L\negthinspace G$.  The
bundle over $S^2$ that is used in constructing the 't Hooft
operator $T(^L\negthinspace R)$ is topologically non-trivial if
and only if the center of $^L\negthinspace G$ acts nontrivially on
$^L\negthinspace R$.

\par\medskip{\it Combined Wilson-'t Hooft Operators: Abelian Case}

We have so far found Wilson line operators at $\Psi=\infty$ and 't
Hooft line operators at $\Psi=0$.  By applying
 duality transformations, we find other special values of
$\Psi$ at which line operators exist that preserve the topological
symmetry.  In fact, all rational values of $\Psi$, and only those,
can be obtained from $\Psi=0$ (or $\Psi=\infty$) by a duality
transformation.  This is true regardless of whether the duality
group $\Gamma$ is $SL(2,\Bbb{Z})$, which it is for $G$
simply-laced, or one of the Hecke groups that arise when $G$ is
not simply-laced.

To get to rational values of $\Psi$, we will take $t=1$ with
$\theta=2\pi \Psi$. Thus, to the action of the gauge field, we add
a term % 式 (6.10)
\setcounter{equation}{9}
\begin{equation}\label{eq:into}
{{\cmmib I}}_\theta=i{\Psi\over 4\pi}\int
\Tr\,F\wedge F.
\end{equation}
 At this stage, it is necessary to assume that $M$
is oriented, say by the choice of a volume-form $\epsilon_4$.
 ${\EUN}=4$ super Yang-Mills theory makes sense on an
unorientable four-manifold $M$.  But the interaction \eqref{eq:into}, whose
role we now wish to consider, can only be introduced if $M$ is
oriented, since this is required in order to define the integral
$\int_M\Tr\,F\wedge F$.  (More generally, $\Psi$ is odd under
orientation-reversal, as we noted at the end of section 第~\ref{sec:canonpar}~节,
so to get to nonzero rational $\Psi$, we require an orientation on
$M$.) So we assume that $M$ is oriented, and express the
orientation by the choice of a volume-form $\epsilon_4$.

The definition of a Wilson loop operator supported on ${\cal S}$
required an orientation of ${\cal S}$ by a one-form $ds$.  The
definition of an 't Hooft loop operator supported on ${\cal S}$
required an orientation $\epsilon_3$ of the normal bundle to
${\cal S}$.  Once $M$ is oriented, we can ask for the orientations
of ${\cal S}$ and its normal bundle to be compatible in the sense
that % 式 (6.11)
\setcounter{equation}{10}
\begin{equation}\label{eq:delo}
\epsilon_4=ds\wedge \epsilon_3
\end{equation}
 along ${\cal S}$.
Identifying a small neighborhood of ${\cal S}$ in $M$ with ${\cal
S}\times \Bbb{R}^3$, we can extend $ds$ over this neighborhood so
that \eqref{eq:delo}\ remains true throughout ${\cal S}\times \Bbb{R}^3$.
Then writing $\epsilon_3=dr\wedge d{\rm Vol}$ as in \eqref{eq:congo}, we
have % 式 (6.12)
\setcounter{equation}{11}
\begin{equation}\label{eq:nelo}
\epsilon_4=ds\wedge dr\wedge d{\rm Vol}=-dr\wedge
ds\wedge d{\rm Vol}.
\end{equation}

The line operators that preserve the topological symmetry at
rational values of $\Psi$ are called mixed Wilson-'t Hooft
operators.  Because these operators will combine the properties of
't Hooft and Wilson operators, to define them it is important to
choose compatible orientations on ${\cal S}$ and its normal
bundle, as we did in the last paragraph. To construct these
operators, we first consider the case $G=U(1)$.

At first sight it may appear that the 't Hooft operators as we
have defined them already possess the topological symmetry at any
$\Psi$.  This is in fact not true, because the singularity of the
't Hooft operator causes ${{\cmmib I}}_\theta$ not to be invariant
under the topological symmetry.

This requires some explanation.  We let $A_0$ be a fixed
connection whose curvature has the Dirac monopole singularity
along a curve ${{\cal S}}$.  A path integral in the presence of a
charge $m$ operator $T_m({{\cal S}})$ is evaluated as follows. One
takes the gauge field to be $A=mA_0+\hat A$, where $\hat A$ is
smooth near ${{\cal S}}$. The path integral is then evaluated by
integrating over $\hat A$, with $A_0$ fixed.  The precise choice
of $A_0$ does not matter (as long as it has the right
singularity), since we are going to integrate over all $\hat A$.

As a step in this direction, let us evaluate ${{\cmmib I}}_\theta$
as a function of $A$.  We write $F=mF_0+\hat F$, where $F_0$ and
$\hat F$ are, respectively, the curvatures of $A_0$ and $\hat A$.
Since we are going to encounter some singularities, we introduce a
cutoff. Letting $V_\epsilon$ be a tube of radius $\epsilon$
centered on ${{\cal S}}$, we will evaluate $\int_{M-V_\epsilon}
F\wedge F$ and consider the limit as $\epsilon\to 0$.

One term comes from $\int_{M-V_\epsilon}m^2 F_0\wedge F_0$.  This
term is not very interesting, because it is independent of the
integration variable $\hat A$.\footnote{This contribution to ${\cmmib
I}_\theta$ actually converges as $\epsilon\to 0$, because the
singular part of $F_0$ is of rank two.} Another contribution that
is not interesting for our present purposes is
$\int_{M-V_\epsilon}\hat F\wedge \hat F$.  Since $\hat A$ is
smooth, this has a limit as $\epsilon\to 0$, namely $\int_M \hat
F\wedge \hat F$. The limiting functional is a topological
invariant, and so is certainly invariant under the change in $\hat
A$ generated by the topological symmetry.

The interesting contribution is the cross term $\tilde {\cmmib
I}_\theta =(im\Psi/2\pi)\int_{M-V_\epsilon}\,F_0\wedge \hat F$. We
want to determine if this expression is annihilated by the
generator $\delta_T$ of the topological symmetry. We define the
action of $\delta_T$ on $A_0$ and $\hat A$ by $\delta_T A_0=0 $
(since $A_0$ is a fixed gauge field) and $\delta_T\hat
A=\delta_TA$. We have $\delta_T \hat F=d (\delta_T\hat A)$, so
% 式 (6.13)
\setcounter{equation}{12}
\begin{equation}\label{eq:juqu}
\delta_T\tilde {{\cmmib I}}_\theta={im\Psi\over
2\pi}\int_{M-V_{\epsilon}}d\left(F_0\wedge \delta_T\hat
A\right)={im\Psi\over 2\pi}\int_{\partial V_\epsilon} F_0\wedge
\delta_T\hat A.
\end{equation}
 To justify the integration by parts here, we must
orient $\partial V_\epsilon$ by contracting the outward normal
vector to $M-V_\epsilon$ with the orientation form of $M$. As the
outward normal to $M-V_\epsilon$ is $-d/dr$, we see from \eqref{eq:nelo}\
that $\partial V_\epsilon$ must be oriented by $ds\wedge d{\rm
Vol}$.

Now the boundary of $V_\epsilon$ is a two-sphere bundle over the
curve ${{\cal S}}$.  The integral over the fibers of
$V_\epsilon\to {{\cal S}}$ can be performed using the fact that
the integral of $F_0$ over a fiber is $-2\pi i$.   So we get
% 式 (6.14)
\setcounter{equation}{13}
\begin{equation}\label{eq:hugu}
\delta_T {{\cmmib I}}_\theta=m\Psi \int_{{\cal
S}}\delta_T \hat A.
\end{equation}

So ${{\cmmib I}}_\theta$ breaks the topological symmetry in the
presence of the 't Hooft operator, or equivalently, the 't Hooft
operator spoils the topological symmetry when $\theta\not= 0$. How
can we restore the symmetry? Let us remember what one does with
the action ${{\cmmib I}}$ in quantum mechanics: one integrates
over all fields with a factor of $\exp(-{{\cmmib I}})$.\footnote{In
Lorentz signature, this factor would be $\exp(i{{\cmmib I}})$, and
${{{\cmmib I}}}_\theta$ would have an additional factor of $i$,
leading to the same result.}  So it is useful to re-express \eqref{eq:hugu}\
as a formula for the variation of $\exp(-{{\cmmib I}}_\theta)$:
% 式 (6.15)
\setcounter{equation}{14}
\begin{equation}\label{eq:jugu}
\delta_T(\exp(-{{\cmmib
I}}_\theta))=-m\Psi\,\exp(-{{\cmmib I}}_\theta)\, \int_{{\cal S}}
\delta_T\hat A.
\end{equation}

{}From here, it is a short step to see that
% 式 (6.16)
\setcounter{equation}{15}
\begin{equation}\label{eq:flugu}
\delta_T\left(\exp\left(m\Psi\int_{{\cal S}} \hat
A\right)\exp\left(-{{\cmmib I}}_\theta\right)\right)=0.
\end{equation}
 So we can
restore the topological symmetry if we include a Wilson operator
$\exp(m\Psi\int_{{\cal S}} \hat A)$ as an additional factor in the
path integral.   (Notice that we have to use here a naive Wilson
operator, not a supersymmetric one with a contribution from
$\phi$.)

However, the expression $\exp(-n\int_{{\cal S}} \hat A)$ is not
gauge-invariant for generic real $n$.  It is gauge-invariant if
and only if $n$ is an integer. So $n=-m\Psi$ must be an integer if
we are to restore the topological symmetry. To put it differently,
$\Psi$ must be a rational number  $-n/m$. For rational
$\Psi=-n/m$, we have learned that the topological symmetry is
preserved by the product of a charge $n$ Wilson operator
$\exp(-n\int_S \hat A)$ and a charge $m$ 't Hooft operator.
 The product
is called a Wilson-'t Hooft operator. We will denote it as
$WT_{n,m}({{\cal S}})$.

For rational $\Psi$, we write $\Psi=-n/m$ with coprime integers
$n$ and $m$.  Under a duality transformation $\Psi\to
(a\Psi+b)/(c\Psi+d)$, $n$ and $m$ transform by
% [WARN] 公式 \eqn\eqref{eq:yto} 无编号或重复
\begin{equation*}
\tag*{(6.17)}
\left(\begin{matrix} m  & n \\\end{matrix}\right)\to  \left(\begin{matrix} m &
n\\\end{matrix}\right) \left(\begin{matrix} d & -b\\ -c & a \\\end{matrix}\right) .
\end{equation*}
 As in
eqn. \eqref{eq:xilj}, this is the standard transformation of the magnetic
and electric charges $\eurm m$ and $\eurm n$ under duality.  The
reason for this is that the state obtained by quantizing the
abelian gauge theory in the presence of a static Wilson-'t Hooft
operator $WT_{n,m}$ (and no other charges or sources) has $({\eurm
m}~{\eurm n})=(m~n)$.  For the magnetic charges, this is fairly
clear from the fact that $m$ and ${\eurm m}$ have both been
defined in terms of the first Chern class (measured either at
infinity or near the singularity). For the electric charges, it
can be demonstrated by solving Maxwell's equations  to determine
the asymptotic behavior of the electric field in the presence of
the Wilson operator.  The computation we have just performed shows
that $n=-\Psi m$ transforms under $\Psi\to\Psi+1$ as claimed in
\eqref{eq:yto}.

\par\medskip{\it Combined Wilson-'t Hooft Operators:
Nonabelian Case}

Just as for 't Hooft operators, the generalization of Wilson-'t
Hooft operators to a general simple gauge group $G$ is made by
embedding the abelian construction in $G$, using a homomorphism
$\rho:U(1)\to G$.  This, however, requires some explanation.

We  write $A=\rho(A_0)+\hat A$, where $A_0$ is the singular $U(1)$
gauge field with the Dirac singularity along ${{\cal S}}$, and
$\hat A$ is smooth.  (For $G=U(1)$, the homomorphism $\rho$ is at
the Lie algebra level multiplication by $m$ for some integer $m$,
so we wrote $A=mA_0+\hat A$.)  Now, however, we must clarify what
it means for $\hat A$ to be smooth.  Because of the nonlinear term
in the curvature $F=dA+A\wedge A$, we must require
% 式 (6.18)
\setcounter{equation}{17}
\begin{equation}\label{eq:milfo}
[\rho(A_0),\hat A]=0
\end{equation}
 along ${{\cal S}}$, or else the
singular part of the curvature along ${{\cal S}}$ will depend on
$\hat A$ and will not coincide with the Dirac singularity.

For the condition \eqref{eq:milfo}\ to make sense, we must also restrict the
gauge transformations along ${{\cal S}}$ to the subgroup $H\subset
G$ that commutes with $\rho(U(1))$.  So, when restricted to
${{\cal S}}$, $\hat A$ has structure group $H$.

The Lie algebra of $H$ has an $H$-invariant projection $\pi_\rho$
to $\rho(\frak{u}(1))$, the Lie algebra of $\rho(U(1))$. Along
${{\cal S}}$, we can define the projection $\pi_\rho(\hat A)$, a
$\rho(\frak{u}(1))$-valued gauge field.

The same calculation as in the abelian case shows that the 't
Hooft operator that corresponds to $\rho$ is no longer invariant
under the topological symmetry  if $\Psi\not= 0$. Instead,
% 式 (6.19)
\setcounter{equation}{18}
\begin{equation}\label{eq:yugo}
\delta_T {{\cmmib I}}_\theta={\Psi}\int_{{\cal S}}
\Tr\,\rho(1)\,\delta_T\hat A.
\end{equation}
 Here $\rho(1)$ is simply the image
of $1\in \frak u(1)$ under the Lie algebra homomorphism
$\rho:\frak u(1)\to \frak g$.  The expression $\Tr\,\rho(1)\,
\delta_T \hat A$ is a multiple of $\pi_\rho(\delta_T\hat A)$.

Just as in the abelian case, topological invariance can be
restored if we multiply the 't Hooft operator by
 % 式 (6.20)
\setcounter{equation}{19}
\begin{equation}\label{eq:yrot}
\exp\left(\Psi\int_{{\cal S}} \Tr \,\rho(1)\hat A\right).
\end{equation}
 For this to be
gauge-invariant places a condition on $\Psi$, which informally is
that $\rho$ must be divisible by the denominator of $\Psi$.  For
any rational $\Psi$, this condition is obeyed for suitable $\rho$.

The operators obtained this way are called Wilson-'t Hooft
operators.  The underlying supersymmetric gauge theory has more
general Wilson-'t Hooft operators with linearly independent
electric and magnetic weights \cite{ref90}.  The Wilson-'t Hooft
operators in the topological field theory all arise by duality
from Wilson operators at $\Psi=\infty$, so their electric and
magnetic weights are proportional.

See also \cite{ref92} for a recent related
discussion of Wilson-'t Hooft operators.
